\section{Introduction}
\label{sec:intro}

A central question for unifying accounts of neural representations is
what, exactly, converges: when systems trained on different problems
arrive at similar internal structure, which property of the problems
predicts the shared solution? We study a setting where the answer can be
computed in closed form before training. The word problem of a finite
group $G$, mapping a sequence of generators to its composed product,
admits an exact algebraic description of the minimal structure any
faithful solution must carry: $\dmin(G)$, the dimension of $G$'s smallest
faithful real linear representation. The loss target already fixes
$\dmin$ as the informative rank; the open question is whether the model's
own recruited capacity, two ambient dimensions wider than $\dmin$,
settles there too, tracking the group's algebra, or instead sorts with
the group's computational class,
solvable versus non-solvable, the divide at which word problems become
$\mathrm{NC}^1$-complete \citep{barrington1989} and at which
state-tracking expressivity results separate recurrent architectures
\citep{merrill2024illusion, grazzi2025negative}.

We train one architecture, a Transformer-encoded single-matrix-state
model under a hard bottleneck, on five groups spanning that divide, and
measure the learned state's restricted effective rank against $\dmin$.
Three results. First, convergence to the algebraic minimum: per-group
mean rank lands within 11\% of $\dmin$, all 19 seeds inside a
pre-registered band, with Spearman $\rho = 0.9747$, the design's
tie-capped maximum (Section~\ref{sec:convergence}).
% <!-- evidence: U1 -->
Second, the designed
head-to-head at matched dimension: $S_4$ (solvable) and $A_5$
(non-solvable) share $\dmin = 3$, and a pre-registered equivalence test
declares their learned ranks statistically equivalent, so the pair
coincides on dimension rather than separating on complexity class
(Section~\ref{sec:equivalence}).
% <!-- evidence: U2 -->
Third, the converged
dimension is causally load-bearing: capping rank at $\dmin{-}1$ pins
recovery below the readout's own geometric ceiling in every group, and
at $\dmin$ that ceiling lifts and recovery empirically returns
(Section~\ref{sec:causal}).
% <!-- evidence: U3 -->
A measurement lesson runs
through all three: two instrument defects each produced a plausible
wrong reading before being caught against raw run artifacts
(Appendix~\ref{app:instrument}), and the convergence law above is what
survived the repaired instrument.
